\documentclass[12pt]{article}

% ─────────────────────────────────────────────────────────────────────────────
% Symbolic Substrates (Paper 2) — SKELETON.
% Unifies two former seeds (rmw3 decision 2026-06-26: "P2 absorbe a P_VSL"):
%   - the LANGUAGE lens (agent-as-compiler, grammar, Shannon)  ← from vsl_language
%   - the SUBSTRATE lens (does meaning survive the crossing?)  ← from symbolic_substrates
% Data + provenance: paper_data.vsm (this directory). Empirical spine = EXP-PHI-01/01-bis/02.
% Prose marked [TO WRITE] is for a later Opus writing session; the skeleton compiles.
% ─────────────────────────────────────────────────────────────────────────────

\usepackage[utf8]{inputenc}
\usepackage[T1]{fontenc}
\usepackage{lmodern}
\usepackage[a4paper, margin=2.5cm]{geometry}
\usepackage{setspace}
\onehalfspacing
\usepackage{amsmath, amssymb}
\usepackage{booktabs}
\usepackage{listings}
\usepackage{xcolor}
\usepackage{microtype}
\usepackage{hyperref}
\hypersetup{colorlinks=true, linkcolor=black, citecolor=teal, urlcolor=teal}

% VSL code style
\lstdefinelanguage{VSL}{
  keywords={V, S1, S2, S3, S4, S5, viable, blocked, partial, true, false,
            recursive, phi_up, phi_down, algedonic, computed, interpreted,
            channel, status, domain_context, env_sample, predictive_alert, veto},
  sensitive=true,
  comment=[l]{//},
  morestring=[b]",
}
\lstset{
  language=VSL,
  basicstyle=\ttfamily\small,
  keywordstyle=\bfseries,
  commentstyle=\itshape\color{gray},
  frame=single,
  breaklines=true,
  columns=fullflexible,
}

\newcommand{\R}{\mathbb{R}}
\newcommand{\nucleus}{\texttt{vsf-nucleus}}
\newcommand{\phiup}{\ensuremath{\varphi\!\uparrow}}
\newcommand{\phidown}{\ensuremath{\varphi\!\downarrow}}
\newcommand{\TODO}[1]{\medskip\noindent\textcolor{gray}{\small[\textsc{to write} --- #1]}\medskip}

\title{Symbolic Substrates:\\
  \large VSL as a Viable Coordination Medium\\
  Between Heterogeneous Cognitive Systems}

\author{Luis Pineda R.\thanks{Contact: \texttt{luis@tallerpineda.cl}} \\
  \small\textit{Viable System Framework Project}}

\date{\textit{Draft skeleton --- June 2026}}

\begin{document}
\maketitle

% ─────────────────────────────────────────────────────────────────────────────
\begin{abstract}
We describe VSL (Viable System Language), a domain-specific language derived from
Beer's Viable System Model, and ask of it two questions that turn out to be one. The
first is what VSL \emph{is}: a notation simultaneously parseable by a deterministic
grammar, generatable by a cognitive agent (a large language model), and semantically
precise enough to govern a multi-layer operational system. We argue that an agent
generating VSL is not translating but \emph{compressing} --- extracting operational
intent from natural language and encoding it in the minimum formal representation that
preserves that intent. This \emph{agent-as-compiler} framing makes language quality
measurable: parse rate, semantic fidelity, round-trip preservation.

The second question is whether VSL \emph{works} as a coordination substrate between two
\emph{different} minds. Here a parse guarantee is necessary but not sufficient: when one
agent encodes intent and a heterogeneous agent decodes it, a correct parse tree with the
wrong referents is a silent failure. We measure this directly and find a sharp split ---
\emph{structural} fidelity crosses at $\approx 100\%$ while \emph{semantic} fidelity
drifts to $\approx 40\%$, the receiving model substituting its own operational context
for the sender's. A lightweight domain anchor ($\approx 8$ tokens) attenuates the drift;
but the load-bearing result is architectural. By partitioning the message into
\emph{computed} fields (executed deterministically by a wrapper) and \emph{interpreted}
fields (executed by the agent), fidelity on the computed axis becomes $1.0$ by
construction --- Ashby's Law of Requisite Variety applied at the level of the grammar,
with computational variety routed to a deterministic layer and semantic variety to the
cognitive one. Meaning survives the crossing not because the model is large enough, but
because the substrate stops asking the model to do what only a parser can.

We report compression ($38$--$79\%$ against natural language, growing with protocol
complexity), the structural/semantic drift and its anchor, the computed/interpreted
split (computed fidelity $=1.0$ across five tasks; $18\times$ message compression), a
capability taxonomy whose boundary is domain knowledge rather than task type, and the
substrate's measurable \emph{cost of existing}. We close with the information-theoretic
reading: the grammar is a codebook and the set of valid VSL files is a low-dimensional
manifold in string space, so generating VSL is finding a point on that manifold.
\end{abstract}

\textbf{Keywords:} viable system model; domain-specific language; LLM-as-compiler;
semantic compression; requisite variety; second-order cybernetics; coordination
substrate; Shannon entropy; VSL.

\medskip\hrule\medskip
\tableofcontents
\newpage

% ─────────────────────────────────────────────────────────────────────────────
\section{The Problem of Saying What a System Is}
\label{sec:problem}
% ─────────────────────────────────────────────────────────────────────────────

% NARRATIVE — origin story. The first handoff too long for a context window but too
% structured for prose. Why NL / JSON / code / math each fail in a specific way.
% VSL emerged bottom-up from the .vsm files we were already writing.

\TODO{origin story: the first session handoff that was too long for a context window
but too structured for prose; the gap was not accidental.}

Four candidate media each fail in a specific way: natural language is ambiguous and
verbose (ten to fifty times the tokens of the state it describes, with no parse
guarantee); JSON/YAML parse but are semantically flat (no notion of system layer,
channel type, or viability constraint); general-purpose code is Turing-complete but
domain-blind (a struct does not know it is an S3 node); and mathematics is precise but
not generatable in real-time interaction. The S2 coordination channel between cognitive
agents had no adequate medium.

\TODO{develop each failure mode; introduce the deeper, second-order problem: even a
medium that parses may not carry meaning across two different minds.}

% ─────────────────────────────────────────────────────────────────────────────
\section{A Grammar Derived From Structure, and the Agent That Compiles It}
\label{sec:language}
% ─────────────────────────────────────────────────────────────────────────────

VSL expresses three things and deliberately omits the rest: system identity
$\mathtt{V}(\mathit{id}, L)$, channel contracts (\phiup, \phidown, algedonic), and
viability state (\texttt{viable}, \texttt{blocked}, \texttt{partial}). Its recursion
mirrors the VSM's: any S1 unit may embed a full system.

\begin{lstlisting}
@vsm 0.1
V(vsf-nucleus, L1) {
  status:  viable
  channel: phi_up    -> V(vsf-s3-star, L2)
  channel: phi_down  <- V(vsf-s3, L2)
  channel: algedonic -> V(vsf-s5, L3)
  recursive: false
}
\end{lstlisting}

\TODO{the BNF fits a page; every production has a reason; what VSL omits (implementation
detail, code, prose) is a design decision, not a limitation.}

\subsection{The agent as compiler}
\label{sec:agent}

A classical compiler maps source to machine code: deterministic, lower abstraction, no
ambiguity resolution, correctness binary. An agent generating VSL maps natural language
(high ambiguity, low precision, high expressiveness) to VSL (low ambiguity, high
precision, constrained expressiveness). This is not classical compilation but
\emph{semantic compression}: the agent extracts operational intent and encodes it in
the minimum formal representation that preserves it. Its quality is therefore a
distribution --- parse rate, semantic fidelity, round-trip preservation --- not a bit.
The constraint that makes this clean is D6 (no language model at L0/L1): the agent
operates at $L \ge 2$, generating VSL that governs L0/L1 processes. The
cognitive/operational boundary \emph{is} the language boundary. In our architecture
Qwen acts as S1, Claude as S3--S5, and VSL is the S2 channel between them.

\TODO{develop agent-as-compiler vs classical compiler; the measurable-quality claim sets
up the experiments of Sections~\ref{sec:substrate}--\ref{sec:measuring}.}

% ─────────────────────────────────────────────────────────────────────────────
\section{Does Meaning Survive the Crossing?}
\label{sec:substrate}
% ─────────────────────────────────────────────────────────────────────────────

A parse guarantee is necessary but not sufficient. When agent $A$ encodes intent as VSL
and a \emph{heterogeneous} agent $B$ decodes it, the question is whether $B$ reconstructs
$A$'s operational referents or its own. A correct parse tree with the wrong referents is
a \emph{silent} failure --- the second-order problem (the cybernetics of observing
systems) made concrete.

We measured it. Submitting schema-mode VSL to an independent model with the VSF schema
preloaded, structural fidelity --- correct S-layer hierarchy, \phiup/\phidown\ channels,
PSI thresholds, the covenant step, metric names --- was $\approx 100\%$. Semantic
fidelity was $\approx 40\%$: the model systematically substituted its own context for
VSF referents (\phiup\ checkpoints mapped to cloud-sync; \texttt{vm-nucleus} to an
internal knowledge module). The single notion of ``fidelity'' had conflated two
independent properties.

\begin{table}[!htb]
\centering
\begin{tabular}{lcc}
\toprule
Property & Crosses the substrate at & Failure mode \\
\midrule
Structural fidelity & $\approx 100\%$ & --- \\
Semantic fidelity (no anchor)   & $\approx 40\%$ & context substitution (silent) \\
Semantic fidelity (with anchor) & $\approx 95\%$ & --- \\
\bottomrule
\end{tabular}
\caption{Structural vs.\ semantic fidelity across a heterogeneous crossing
(EXP-PHI-01). Structure is reconstructed almost perfectly; meaning drifts unless
anchored. A \texttt{domain\_context} header of $\approx 8$ tokens attenuates the drift
from $\approx 60\%$ to $\approx 5\%$.}
\label{tab:drift}
\end{table}

\begin{lstlisting}
domain_context: {target=vsf-nucleus, env=Rust/LXC, layer=S1}
\end{lstlisting}

\TODO{develop context substitution as a structural property of heterogeneous coupling,
not a model defect; the anchor as a cheap variety attenuator; set up the deeper fix.}

% ─────────────────────────────────────────────────────────────────────────────
\section{Requisite Variety at the Grammar Level}
\label{sec:ashby}
% ─────────────────────────────────────────────────────────────────────────────

The anchor attenuates the drift; it does not eliminate the cause. The cause is that we
were asking the cognitive layer to perform computational work --- counting, matching,
enumerating --- at which a parser is exact and a language model is merely probable. The
fix is Ashby's: route variety to the layer competent to handle it. We partition the
message's \texttt{output\_format} into \emph{computed} fields, executed deterministically
by a wrapper before the agent is ever called, and \emph{interpreted} fields, executed by
the agent alone.

\begin{lstlisting}
output_format: {
  computed:    { line_count: INTEGER, token_frequencies: OBJECT, regex_matches: ARRAY }
  interpreted: { summary: STRING, key_themes: ARRAY, cybersyn_mapping: OBJECT }
}
\end{lstlisting}

On the computed axis, fidelity is $1.0$ by construction: there is no model variance to
contend with, because no model is consulted. On the interpreted axis, fidelity remains
capability-dependent --- which is correct, because that is where genuine interpretation
lives.

\begin{table}[!htb]
\centering
\begin{tabular}{lcc}
\toprule
Axis & Fidelity & Mechanism \\
\midrule
Computed (wrapper)     & $1.00$ (5/5 tasks, 3 runs) & deterministic by design \\
Interpreted (agent)    & capability-dependent        & domain knowledge of the model \\
\midrule
Message size (split)   & $\approx 450$ B vs.\ $8000$ B & $18\times$ compression \\
\bottomrule
\end{tabular}
\caption{The computed/interpreted split (EXP-PHI-02). Partitioning by competence makes
computed fidelity a design invariant rather than a model outcome, architecturally
resolving the closure-is-not-correctness problem of EXP-PHI-01. The latency target was
not met under a 1.5\,B model on CPU --- a hardware/model bound, not a property of the
split.}
\label{tab:split}
\end{table}

This architecturally closes the earlier finding that \emph{loop closure is not
correctness}: the failure was structural, not a deficit of model scale. A capability
taxonomy follows naturally --- and its boundary is instructive. Reliability on
interpreted fields tracks the model's \emph{domain knowledge}, not the task type: a
1.5\,B model correctly summarised an ethics specification but read ``covenant'' as a
legal contract, and mapped a governance \texttt{watchdog} to S3 (control) rather than S5
(policy). The boundary that matters is not counting-vs-reasoning but in-domain-vs-out.

\TODO{develop the Ashby framing as the paper's central constructive result; the
capability matrix; the two boundaries (task-type robust, domain-knowledge model-specific).}

% ─────────────────────────────────────────────────────────────────────────────
\section{The Cost of Depending on Another Mind}
\label{sec:sovereignty}
% ─────────────────────────────────────────────────────────────────────────────

% H5 substrate_sovereignty: the multi-dimensional cost of existing.

A substrate that depends on an external cognitive agent is not free. Its cost of existing
is multi-dimensional --- credentials, routing, upstream availability, dynamic resources
--- and it is structural, surfacing in the loop rather than as an accident of any one
run. This connects the substrate to the broader cost-of-coherence thread.

\TODO{develop H5; the sovereignty cost as a measurable property; the bridge to the
epistemic-viability companion paper.}

% ─────────────────────────────────────────────────────────────────────────────
\section{Measuring the Language}
\label{sec:measuring}
% ─────────────────────────────────────────────────────────────────────────────

\subsection{Compression}
\label{sec:compression}

% EXP-PHI-01 compression table (tiktoken cl100k_base).

\begin{table}[!htb]
\centering
\begin{tabular}{lrrrrr}
\toprule
Case & $T_A$ & $T_{B,\text{full}}$ & $T_{B,\text{schema}}$ & $r_\text{full}$ & $r_\text{schema}$ \\
\midrule
\texttt{pvsnp}               & 328 & 161 & ---  & 0.51 & ---  \\
\texttt{thermal\_senescence} & 638 & 346 & ---  & 0.46 & ---  \\
\texttt{spec\_boot+Lyapunov} & 717 & 432 & 148  & 0.40 & 0.79 \\
\texttt{proxmox+S4}          & 698 & 411 & 198  & 0.41 & 0.72 \\
\texttt{PSI/1.0}             & 641 & 397 & 147  & 0.38 & 0.77 \\
\bottomrule
\end{tabular}
\caption{VSL compression (tiktoken \texttt{cl100k\_base}). Standalone VSL achieves
$38$--$51\%$; with a shared schema (machine-to-machine), $72$--$79\%$, growing with
protocol complexity. $r$ is the compressed-to-original token ratio.}
\label{tab:compression}
\end{table}

\subsection{Parse rate and generation}
\label{sec:parse}

The VSL v1.0 parser accepts $212/217$ corpus files ($\approx 97.7\%$); the five failures
are legacy-format edge cases, not grammar violations. An S1 agent (Qwen) generates valid
\texttt{next\_experiment.vsm} documents through a plan/eval/propose loop, the operational
demonstration that the grammar is learnable from the specification alone.

\subsection{Shannon entropy of VSL versus natural language}
\label{sec:entropy}

% EXP-VSL-04 — run 2026-06-29; results in docs/experiments/EXP-VSL-04/results.vsm

EXP-VSL-04 measured Shannon entropy over five paired specifications ---
each written once in natural-language prose and once in VSL with equivalent
operational content --- at two granularities: character-level, and
token-level under a production tokenizer (\texttt{cl100k\_base}).
The two levels dissociate, and the dissociation is the finding.
At the character level VSL carries \emph{more} entropy than prose
($H_\text{char}$: $5.52$ vs.\ $4.99$ bits, $+10.6\%$): the notation's symbol
inventory (\texttt{@}, \texttt{:=}, braces, arrows) widens the character
distribution. At the token level the order reverses
($H_\text{tok}$: $6.89$ vs.\ $7.28$ bits, $-5.4\%$, holding in all five
pairs): the grammar constrains how symbols combine, so the token stream is
more predictable than prose even though the character stream is less so.
The codebook operates at the token level --- the level at which an agent
consumes the specification. Combined with the token-count compression of
Table~\ref{tab:compression} (38--79\% fewer tokens for the same content),
lower per-token entropy over fewer tokens is the information-theoretic
signature of a constrained grammar doing real coding work.

Two boundaries on the claim. First, this is evidence of a \emph{good} code,
not an optimal one: optimality against the source entropy of ``viable-system
specifications'' would require a rate floor we cannot estimate, and
``semantic entropy'' has no agreed definition we are willing to build on.
The operational proxy we adopt instead is task-conditional: a code is good
to the extent that a competent agent reconstructs the specified system from
fewer tokens at equal fidelity --- Section~\ref{sec:substrate}'s experiment
restated in coding-theoretic terms. Second, the schema-only variants
compress further but raise character entropy further ($+0.18$ to $+0.67$
bits over prose): stripping inline definitions concentrates notation, a
rate--redundancy trade the grammar makes visible rather than hides.

% ─────────────────────────────────────────────────────────────────────────────
\section{Grammar as a Mirror of the Architecture}
\label{sec:mirror}
% ─────────────────────────────────────────────────────────────────────────────

Every gap in VSL v1 corresponds to an architectural debt; VSL v2 closes those gaps in
parallel with the system closing them. A rejected-\phidown\ channel appears once S1
validates \phidown\ before applying it; a \texttt{predictive\_alert} event and an
\texttt{env\_sample} channel appear once the system has an external environment to sense
(the companion kernel paper shows the internal anticipatory channel is empty without
external variance, motivating exactly this grammar); a \texttt{veto} channel, distinct
from the algedonic alert, appears once S5 can halt as well as warn. The grammar does not
run ahead of the architecture --- it is precisely aligned with it, and its honesty about
what it cannot yet express is part of the contribution.

The mirror thesis yields one forward-looking, falsifiable prediction. The
architecture is recursive --- any S1 unit may itself be a complete VSM ---
and if the grammar truly mirrors the architecture, the \emph{specification}
should inherit that self-similarity: the compression ratio
$r_L = |V_L|/|V_{L-1}|$ between a level-$L$ spec and the specs of its
level-$(L\!-\!1)$ constituents should be approximately constant across
recursion levels. A two-level recursive deployment (planned) tests this
directly. We state it as a prediction about description length only; we
make no claim that this ratio relates to the dynamical contraction rate of
the kernels themselves, which lives in a different space.

\TODO{develop grammar-as-architecture-mirror; v2 as the documented co-evolution; keep
the claims carried by v1 (H1/H3/H5/H6/H7/H8/H9) distinct from v2 designed extensions.}

% ─────────────────────────────────────────────────────────────────────────────
\section{What a Substrate Is For}
\label{sec:closing}
% ─────────────────────────────────────────────────────────────────────────────

% NARRATIVE — closing argument.

A language is successful when it compresses the right things. VSL compresses operational
state; S5 policy deliberately remains in natural language, because policy carries an
intentionality VSL is not designed to encode --- the boundary is a choice, not a defect.
The grammar is a codebook and valid VSL is a low-dimensional manifold in the space of all
strings; the agent that generates VSL is finding a point on that manifold, and the
substrate is viable precisely when the manifold is reachable by a cognitive agent and the
computational variety has been routed elsewhere.

\TODO{closing argument; open questions: emergent language (can an agent learn VSL without
being taught?); can VSL express its own cost-of-existing?}

% ─────────────────────────────────────────────────────────────────────────────
\begin{thebibliography}{9}
\bibitem{beer1972} Beer, S. (1972). \textit{Brain of the Firm}. Allen Lane.
\bibitem{beer1979} Beer, S. (1979). \textit{The Heart of Enterprise}. Wiley.
\bibitem{ashby1956} Ashby, W.~R. (1956). \textit{An Introduction to Cybernetics}. Chapman \& Hall.
\bibitem{shannon1948} Shannon, C.~E. (1948). A Mathematical Theory of Communication.
  \textit{Bell System Technical Journal}, 27.
\bibitem{vonfoerster1974} von Foerster, H. (1974). \textit{Cybernetics of Cybernetics}.
  BCL, University of Illinois.
\bibitem{medina2011} Medina, E. (2011). \textit{Cybernetic Revolutionaries}. MIT Press.
\end{thebibliography}

\end{document}
